\section{Mathematical foundations}
\label{sec:foundations}

\subsection{Markov decision processes and the objective}
\label{sec:mdp}

\begin{definition}[Markov decision process]
A \emph{Markov decision process} (MDP) is a tuple
$\mathcal{M} = (\mathcal{S}, \mathcal{A}, P, r, \gamma)$: a set of states
$\mathcal{S}$; a set of actions $\mathcal{A}$; a transition kernel
$P(s' \mid s,a)$; a reward function $r : \mathcal{S} \times \mathcal{A} \to \R$;
and a discount factor $\gamma \in [0,1]$.
\end{definition}

A \emph{policy} maps states to distributions over actions. When the policy is
represented by parameters $\theta$, we write $\pi_\theta(a \mid s)$. The value
functions are

\begin{align}
v^{\pi}(s) &= \E_{\pi}\!\left[\textstyle\sum_{t \ge 0} \gamma^t r_t
\;\middle|\; s_0 = s\right], &
Q^{\pi}(s,a) &= \E_{\pi}\!\left[\textstyle\sum_{t \ge 0} \gamma^t r_t
\;\middle|\; s_0 = s, a_0 = a\right],
\end{align}

and they satisfy the Bellman equations
$v^\pi(s) = \sum_a \pi(a\mid s) Q^\pi(s,a)$ and
$Q^\pi(s,a) = r(s,a) + \gamma \sum_{s'} P(s'\mid s,a)\, v^\pi(s')$.

The \emph{objective} is the expected outcome of the initial states under the
policy,
\begin{equation}
J(\theta) \;=\; \E_{s_0 \sim \rho}\!\left[ v^{\pi_\theta}(s_0) \right],
\label{eq:objective}
\end{equation}
where $\rho$ is the distribution over starting states.

Two specializations are worth stating at once, because the rest of this
monograph uses them exclusively.

\begin{remark}[Two-player perfect-information games as MDPs]
For a deterministic, perfect-information, zero-sum game with two players, the
following hold.
\begin{enumerate}
  \item $\mathcal{S}$ is the set of legal board positions, $\mathcal{A}$ the set
        of legal moves at a position, and $P$ is deterministic: a move
        determines the next position.
  \item The reward is zero at every non-terminal step and equals the game
        outcome at a terminal state. With outcomes in $\{-1, 0, +1\}$ and
        $\gamma = 1$, the return from a terminal state is exactly the game
        result, and every value lies in $[-1,1]$.
  \item The two players' value functions are related by $v(s) = -v(s')$ when
        $s'$ is the position seen by the opponent. This \emph{sign alternation}
        is the single most common source of implementation defects
        (Section~\ref{sec:tree}).
\end{enumerate}
\end{remark}

\subsection{The likelihood-ratio identity, and REINFORCE}
\label{sec:reinforce}

The objective \eqref{eq:objective} is an expectation over trajectories, and
the trajectories depend on $\theta$. The first step of the whole subject is to
make that dependence disappear by differentiating a probability instead of a
value.

Let $\tau = (s_0, a_0, s_1, a_1, \dots, s_T)$ be a trajectory. Its probability
under the policy factorizes:

\begin{equation}
p_\theta(\tau) \;=\; \rho(s_0) \prod_{t=0}^{T-1}
\pi_\theta(a_t \mid s_t)\, P(s_{t+1} \mid s_t, a_t).
\label{eq:trajprob}
\end{equation}

\begin{proposition}[Likelihood-ratio identity]
\label{prop:likelihood-ratio}
For any function $f$ of the trajectory,
$\nabla_\theta\, \E_{\tau \sim p_\theta}[f(\tau)] =
\E_{\tau \sim p_\theta}\!\left[ f(\tau)\, \nabla_\theta \log p_\theta(\tau) \right]$.
\end{proposition}

\begin{proof}
$\nabla_\theta \E[f] = \nabla_\theta \int p_\theta(\tau) f(\tau) \dd\tau
= \int f(\tau) \nabla_\theta p_\theta(\tau) \dd\tau$. Since
$\nabla_\theta p_\theta = p_\theta \nabla_\theta \log p_\theta$, the identity
follows, provided the exchange of derivative and integral is justified; this
requires that $\pi_\theta$ be differentiable in $\theta$ and that the standard
dominated-convergence conditions hold.
\end{proof}

Apply Proposition~\ref{prop:likelihood-ratio} with $f(\tau) = R(\tau)$, the
return of the trajectory. Taking logarithms in \eqref{eq:trajprob}, the initial
state distribution and the transition kernel do not depend on $\theta$, so all
that remains is

\begin{equation}
\nabla_\theta \E[R] \;=\;
\E\!\left[ R(\tau) \sum_{t=0}^{T-1} \nabla_\theta \log \pi_\theta(a_t \mid s_t)
\right].
\label{eq:reinforce}
\end{equation}

This is the REINFORCE estimator \citep{williams1992reinforce}. It needs no
model of $P$: samples and a score function suffice. With a
\emph{return-to-go} refinement, $R(\tau)$ is replaced at step $t$ by
$G_t = \sum_{t'\ge t} \gamma^{t'-t} r_{t'}$, because rewards before time $t$
cannot be affected by $a_t$; the estimator is unchanged in expectation and
lower in variance.

For a board game with $\gamma = 1$ and outcome $z \in \{-1,0,+1\}$, the
return-to-go from any non-terminal step is simply $z$. Equation
\eqref{eq:reinforce} becomes

\begin{equation}
\nabla_\theta J(\theta) \;=\;
\E\!\left[ z \sum_{t} \nabla_\theta \log \pi_\theta(a_t \mid s_t) \right],
\label{eq:reinforce-game}
\end{equation}

that is: raise the log-probability of every move played in a won game, and
lower it in a lost game. This is exactly the reinforcement-learning stage of
AlphaGo (2016) \citep{silver2016alphago}. It is worth holding on to for two
paragraphs, because it is the point at which the popular description of
AlphaZero as ``policy gradient plus search'' goes wrong
(Section~\ref{sec:loss}).

A \emph{baseline} $b(s)$, any function of the state that does not depend on the
action, leaves the estimator unbiased:

\begin{equation}
\E\!\left[ (G_t - b(s_t))\, \nabla_\theta \log \pi_\theta(a_t \mid s_t) \right]
= \nabla_\theta J(\theta),
\label{eq:baseline}
\end{equation}

because $\E_{a \sim \pi}[b(s) \nabla \log \pi(a|s)] = b(s) \nabla \sum_a \pi(a|s)
= b(s) \nabla 1 = 0$. Choosing $b(s_t) = v^{\pi}(s_t)$ turns the multiplier into
an advantage estimate and reduces variance by orders of magnitude in practice;
choosing $b$ as a bootstrapped estimate is the actor--critic idea
(Section~\ref{sec:actorcritic}).

\subsection{The policy gradient theorem}
\label{sec:pgt}

Equation \eqref{eq:reinforce} has a weakness that is not visible in the board-game
setting but is fatal elsewhere: the multiplier $G_t$ is the return of a whole
sampled trajectory, so its variance grows with the horizon. It would be far
better to weight the score function by a local quantity, $Q^\pi(s_t,a_t)$. The
policy gradient theorem says that this is allowed, and that the price is a
change of state distribution.

\begin{theorem}[Policy gradient theorem, \citeauthor{sutton2000policygradient}]
\label{thm:pgt}
For a differentiable policy $\pi_\theta$ and any start-state distribution,
\begin{equation}
\nabla_\theta J(\theta) \;=\;
\sum_{s} d^{\pi}(s) \sum_{a} Q^{\pi}(s,a)\, \nabla_\theta \pi_\theta(a \mid s),
\label{eq:pgt}
\end{equation}
where $d^{\pi}(s) = \sum_{t\ge 0} \gamma^t \Prob(s_t = s)$ is the
discounted state-visitation distribution under $\pi$.
Equivalently,
$\nabla_\theta J(\theta) = \E_{s \sim d^\pi}
\!\left[\sum_a Q^\pi(s,a) \nabla_\theta \pi_\theta(a|s)\right]$.
\end{theorem}

\begin{proof}[Proof sketch]
Write $J(\theta) = \sum_s d_0(s) v^{\pi_\theta}(s)$ and differentiate
$v^{\pi_\theta}$ by unrolling one Bellman step:
\[
\nabla v^{\pi}(s)
= \sum_a \Big[ \nabla \pi(a|s)\, Q^{\pi}(s,a)
  + \pi(a|s) \sum_{s'} P(s'|s,a)\, \gamma\, \nabla v^{\pi}(s') \Big].
\]
Note which terms are missing: $Q^\pi$ is not differentiated. Substitute this
expression into itself repeatedly. Every unrolling multiplies the state
distribution by $\gamma \sum_a \pi(a|s) P(s'|s,a)$, so after infinitely many
unrollings the outer distribution becomes exactly $d^\pi$, and the term
$\nabla \pi$ is collected once at each level. Summing the resulting geometric
series gives \eqref{eq:pgt}.
\end{proof}

The content of the theorem is what the proof sketch reveals: \emph{the
derivative of the state distribution is never needed}. The distribution
$d^\pi$ appears as a weight, not as a function to be differentiated. This is
what makes sample-based estimation possible at all, and it is why
\eqref{eq:reinforce} and \eqref{eq:pgt} agree in expectation: the sample return
$G_t$ is an unbiased estimate of $Q^\pi(s_t,a_t)$ under the same distribution.

Combining \eqref{eq:pgt} with the log-derivative trick gives the practical form

\begin{equation}
\nabla_\theta J(\theta) \;=\;
\E_{s \sim d^\pi,\, a \sim \pi}\!\left[ Q^{\pi}(s,a)\, \nabla_\theta
\log \pi_\theta(a \mid s) \right],
\label{eq:pgt-score}
\end{equation}

which is the same as \eqref{eq:reinforce} except that the multiplier is now a
\emph{state--action value}, not a trajectory return. Any unbiased estimate of
$Q^\pi$ may be substituted; the choice of estimator is the actor--critic design
space.

\subsection{Actor--critic, bootstrapping, and the deadly triad}
\label{sec:actorcritic}

Substituting a learned approximation $Q_w(s,a) \approx Q^\pi(s,a)$ into
\eqref{eq:pgt-score} gives the actor--critic architecture: an \emph{actor}
$\pi_\theta$ that is improved by the gradient, and a \emph{critic} $Q_w$ that
supplies the multiplier. There are two ways to train the critic, and they
differ in a way that decides the whole stability picture.

\paragraph{Monte Carlo.} Use the observed return $G_t$ as the target for
$Q_w(s_t,a_t)$. The target is an unbiased sample of the true value under the
current policy. The variance is high, because a single game outcome is a noisy
estimate of one position's worth; the bias is zero.

\paragraph{Temporal difference.} Use a one-step bootstrap
$\hat{y}_t = r_t + \gamma\, Q_w(s_{t+1}, a_{t+1})$ as the target, and descend on
the squared temporal-difference (TD) error
\begin{equation}
\delta_t \;=\; r_t + \gamma\, Q_w(s_{t+1},a_{t+1}) - Q_w(s_t,a_t).
\label{eq:td}
\end{equation}
The target is biased, because it depends on the same estimate being trained.
The variance is much lower, because it involves one transition rather than a
whole episode. TD learning propagates information backwards through the state
space, which is why it is sample-efficient; but the bias is real, and it becomes
harmful in combination with the next two ingredients.

\begin{remark}[The deadly triad]
Divergence of value estimation is possible when three conditions hold
simultaneously: \emph{function approximation}, \emph{bootstrapping}, and
\emph{off-policy} learning \citep{suttonbarto2018}. Each is individually
manageable. Approximate value functions alone may simply be inaccurate.
Bootstrapping alone is the basis of TD learning. Off-policy learning alone
requires importance weights but is sound. Together they can produce a feedback
loop in which the approximation error of the target feeds the update, which
enlarges the error.
\end{remark}

This remark returns in Section~\ref{sec:alphazero-why}, because one of the
cleanest features of the AlphaZero design is that it removes the second
ingredient completely: the value target is the raw game outcome, so the loop
never bootstraps.

\subsection{The deterministic policy gradient, and the DDPG branch}
\label{sec:dpg}

The policy gradient of Section~\ref{sec:pgt} integrates over actions. For a
discrete game that integration is cheap: the number of legal moves is at most a
few hundred. For a continuous action space it is the dominant cost, and the
usual resolution is to make the policy deterministic.

\begin{definition}[Deterministic policy]
A deterministic policy is a differentiable map
$\mu_\theta : \mathcal{S} \to \mathcal{A}$. Its objective is
\begin{equation}
J(\theta) \;=\; \E_{s \sim \rho^{\mu}}\!\left[ Q^{\mu}\!\left(s, \mu_\theta(s)\right)\right],
\label{eq:dpg-objective}
\end{equation}
where $\rho^\mu$ is the discounted state distribution induced by $\mu_\theta$.
\end{definition}

Because the action is now a deterministic function of the state, the gradient
of the objective with respect to $\theta$ can be taken by the chain rule -- the
action distribution no longer depends on $\theta$ in a way that must be
integrated out.

\begin{proposition}[Deterministic policy gradient, \citeauthor{silver2014dpg}]
\label{prop:dpg}
Under the usual conditions on differentiability and on the existence of
$\rho^\mu$,
\begin{equation}
\nabla_\theta J(\theta) \;=\;
\E_{s \sim \rho^{\mu}}\!\left[
\underbrace{\nabla_\theta \mu_\theta(s)}_{\text{how to move}}
\cdot
\underbrace{\nabla_a Q^{\mu}(s,a)\big|_{a = \mu_\theta(s)}}_{\text{which way is better}}
\right].
\label{eq:dpg}
\end{equation}
\end{proposition}

\begin{proof}[Proof sketch]
Differentiate \eqref{eq:dpg-objective}. Two kinds of term appear: those in which
$\theta$ enters through $\mu_\theta$ inside $Q^\mu$, and those in which $\theta$
enters through the state distribution $\rho^\mu$. The second kind is
$\sum_s \nabla_\theta \rho^\mu(s)\, Q^\mu(s,\mu_\theta(s))$. Using the
consistency identity of the deterministic policy,
$\nabla_\theta \rho^\mu(s) = \rho(s_0)\sum_t \gamma^t \nabla_\theta P(s_t = s)$,
the state-distribution term can be rewritten as
$\sum_t \gamma^t \sum_s \rho(s_0) \nabla_\theta P(s_t = s)
\, Q^\mu(s,\mu_\theta(s))$. This expression vanishes because it is the
derivative of $\sum_s P(s_t = s) = 1$ with respect to $\theta$, dotted with the
constant factor $Q^\mu(s,\mu_\theta(s))$ evaluated at the same $s$; formally,
it is an expectation of a total derivative of a probability and cancels by the
same argument that removes the state-distribution derivative in
Theorem~\ref{thm:pgt}. What remains is the two-factor product
\eqref{eq:dpg}. The full argument, including the conditions, is
\citet{silver2014dpg}.
\end{proof}

The two factors of \eqref{eq:dpg} have a direct reading that is worth keeping:
$\nabla_\theta \mu_\theta(s)$ says how the chosen action changes when the
parameters change, and $\nabla_a Q^\mu$ says which direction in action space
improves the value. Their product is an ascent direction for the value, obtained
without integrating over actions at all.

The theorem also explains the relationship between the two branches. If the
stochastic policy is a fixed-variance Gaussian,
$\pi_\theta(a|s) = \mathcal{N}\!\left(\mu_\theta(s), \sigma^2 I\right)$, then the
stochastic policy gradient of Theorem~\ref{thm:pgt} converges, as
$\sigma \to 0$, to the deterministic gradient \eqref{eq:dpg} plus a term of order
$\sigma^2$ involving the second derivative of $Q$ at the mean. In the limit, only
the deterministic formula survives. That limit is the content of the connection,
and it is the reason a deterministic actor can be trained at all: the gradient
of the value with respect to the action is a usable teaching signal.

\begin{remark}[Attribution, and a correction worth making once]
The deterministic policy gradient is \citet{silver2014dpg}. The deep version of
it -- an off-policy actor--critic with a replay buffer, target networks, and
action-noise exploration, which is the algorithm usually called \emph{deep
deterministic policy gradient} or DDPG -- is \citet{lillicrap2015ddpg}, whose
author list begins with Lillicrap and Hunt and ends with Silver and Wierstra.
Silver is a co-author of DDPG, not its author, and DDPG is not the parent of
AlphaGo: it is a parallel branch. The two branches share the actor--critic
frame, the chain rule through the action, and nothing else that matters. DDPG
exists because continuous actions make a stochastic policy expensive to
integrate over; AlphaZero exists because discrete games make a lookahead search
cheap to run. DDPG learns a bootstrapped critic $Q_w$ and must inject
exploration noise into a deterministic actor. AlphaZero learns no critic at all:
the search supplies local action values, the value head supplies $v$, and the
target is never bootstrapped. Section~\ref{sec:three-papers} resumes the
AlphaZero branch and does not use DDPG again.
\end{remark}

\subsection{Approximate policy iteration}
\label{sec:apxpi}

Classical policy iteration alternates two steps in an exact MDP:

\begin{enumerate}
  \item \textbf{Policy evaluation}: compute $v^{\pi}$ for the current policy.
  \item \textbf{Policy improvement}: set $\pi'(s) = \argmax_a Q^{\pi}(s,a)$.
\end{enumerate}

The guarantee for the second step is the improvement theorem: acting greedily
with respect to the value of a policy produces a policy that is at least as
good everywhere, and strictly better wherever the greedy action improves on the
policy's own action \citep{suttonbarto2018}. Because the number of deterministic
policies in a finite MDP is finite and the sequence is monotonically improving,
policy iteration converges to an optimal policy in finitely many iterations.

Both steps become approximations as soon as the state space is large. The
improvement step still requires a maximization over actions -- cheap in a
discrete game, impossible in a continuous one -- and the evaluation step
requires a function approximator. The resulting scheme, \emph{approximate policy
iteration}, has no general convergence guarantee, and there are classical
counterexamples.

The AlphaZero loop is an approximate policy iteration, and the AlphaGo Zero
paper states exactly this:

\begin{quote}
``The AlphaGo Zero self-play algorithm can be understood as an approximate
policy iteration scheme in which MCTS is used for both policy improvement and
policy evaluation'' \citep{silver2017alphagozero}.
\end{quote}

The mapping onto the two steps is as follows.

\emph{Improvement} is the tree search. Given the network's policy $p$ and value
$v$, the search returns a distribution $\pi$ over the root's moves. That
distribution is what the network's policy would become if it could look ahead:
it is a better policy than $p$, for the same reasons a greedy step on an
accurate $Q$ is better than the policy that generated the $Q$.

\emph{Evaluation} is the value head trained on $z$, the outcome of games played
with the improved policy. There is no separate critic; the value network is
retrained on outcomes generated by the improved policy, which is precisely the
policy-evaluation step of the loop.

The loop is then: improve with search, evaluate by training on outcomes of the
improved policy, re-derive the policy from the improved network, repeat. The
structural reason it does not stall at the network's own level is the asymmetry
between the two actors: the teacher (the search, which uses the network as a
prior and as a leaf evaluator) is stronger than the student (the network alone),
because the search corrects the network's errors by lookahead and by comparing
alternatives. Every iteration distils that correction back into the weights,
which raises the level at which the next search operates.

\subsection{Expert Iteration}
\label{sec:exit}

The cleanest theoretical frame for the loop is \emph{Expert Iteration} (ExIt),
introduced by \citet{anthony2017exit}. It separates the two roles explicitly.

A \emph{apprentice} is a function approximator $\pi_\theta$ that imitates
decisions. An \emph{expert} is a planner that improves upon the apprentice by
lookahead and returns a better decision distribution. Expert Iteration repeats
one step: plan with the expert on states, then refit the apprentice to the
expert's improved decisions. Formally, with an expert operator
$E(\pi_\theta)$ that maps a policy to a better one,

\begin{equation}
\pi_{\theta} \;\longmapsto\; E(\pi_\theta) \;\longmapsto\;
\argmin_{\theta'} D\!\left(E(\pi_\theta) \, \|\, \pi_{\theta'}\right)
\;\longmapsto\; \pi_{\theta'}, \qquad \text{repeat.}
\label{eq:exit}
\end{equation}

\citet{anthony2017exit} prove that in the exact, tabular case -- an exact expert
and an unrestricted apprentice -- this iteration converges to the optimal policy.
The proof is a polytope argument: the expert's improvement step is a contraction
in a suitable norm, and the imitation step is exact in the tabular case, so the
sequence of policies is monotone.

AlphaGo Zero is Expert Iteration with MCTS as the expert and a residual
convolutional network as the apprentice. Everything the proof needs is lost in
that instantiation: the expert is an approximation (a finite number of
simulations), the apprentice is a deep network with a fixed architecture, the
data distribution moves with the weights, and the imitation target is stochastic.
The value network has no counterpart in ExIt at all; it is an addition that
supplies the expert with a leaf evaluation. The frame is nevertheless the right
one, because it names the two moving parts, and because it identifies exactly
which assumption each practical departure violates. Section~\ref{sec:proven}
returns to this inventory.

\subsection{Bandits, UCT, and PUCT}
\label{sec:bandits}

Tree search in the AlphaZero line is a sequence of bandit problems, one per
node. This is the mathematical content behind the selection rule; the rest of
the search is bookkeeping.

\paragraph{The bandit problem.} In a $K$-armed bandit, a player chooses one of
$K$ options at each of $n$ rounds and observes a bounded reward from the chosen
option. The quality of a strategy is its \emph{regret}: the expected difference
between the reward of always playing the best arm and the reward actually
collected. The problem is the canonical formalization of the
exploration--exploitation trade-off.

\paragraph{UCB1.} \citet{auer2002ucb1} give the rule
\begin{equation}
a_n \;=\; \argmax_{a}\,\left[\, \bar{Q}_a \;+\;
\sqrt{\frac{2 \ln n}{n_a}} \,\right],
\label{eq:ucb1}
\end{equation}
where $\bar{Q}_a$ is the sample mean reward of arm $a$ and $n_a$ its number of
plays. The first term exploits; the second explores, shrinking as $n_a$ grows and
growing as the total number of plays $n$ grows. UCB1 has logarithmic regret per
arm, and the proof is a concentration argument: the exploration bonus is the
width of a confidence interval for $\bar{Q}_a$ at the required confidence level,
and an arm that could still be optimal is never ruled out permanently.

\paragraph{UCT.} \citet{kocsis2006uct} apply the UCB1 rule at every node of a
search tree instead of at a single bandit. A node's "arms" are its actions, the
"reward" of an arm is the value of the state reached by it, and the values are
backed up along the path to the root. The resulting algorithm is Upper Confidence
bounds applied to Trees (UCT). Its theoretical attraction is a consistency
result: as the number of simulations grows without bound, UCT's root decision
converges to the minimax value, with finite-sample error bounds. UCT is the
algorithm that made Monte Carlo Tree Search (MCTS) work in Go, a domain in which
alpha--beta search had failed for decades \citep{coulom2006mcts, browne2012survey}.

\paragraph{PUCT.} The published AlphaGo Zero selection rule adds a learned prior
to the exploration bonus \citep{rosin2011puct}:

\begin{equation}
a^{*} \;=\; \argmax_a \left[\, Q(s,a) \;+\; c_{\mathrm{puct}}\,
P(s,a)\, \frac{\sqrt{\sum_b N(s,b)}}{1 + N(s,a)} \,\right].
\label{eq:puct}
\end{equation}

The three factors of the exploration bonus each do a job. $P(s,a)$ is the prior
supplied by the network's policy head: it concentrates the search where the
network believes the good moves are. $\sqrt{\sum_b N(s,b)}$ makes the bonus grow
as the node is visited more often, so that a node is never abandoned. The
denominator $1 + N(s,a)$ makes the bonus shrink for the specific arm that has
been tried. The sum form of the growth term is the substantive difference from
UCB1: \emph{any} sibling's visit raises the bonus for \emph{this} arm, so
unexplored arms of a heavily visited node keep attracting attention, and the
sequence of unvisited arms is enumerated rather than sampled at random.

The practical difference between UCT and PUCT is the cost of finding out which
arms are good. The classical UCB1 bonus has to spend visits to estimate every
arm's mean, because it starts from total ignorance. In a game with 150 legal
moves, most of which lose immediately or do nothing, that is most of the budget
spent on distinguishing near-equal bad arms. A prior that is even roughly right
focuses the budget on the few moves that matter, and the whole search becomes
deeper instead of wider.

\begin{remark}[What the prior costs, theoretically]
The guarantee does not survive the prior. UCT's consistency proof needs the
exploration term to be a genuine confidence sequence for an agnostic estimator.
Once $P(s,a)$ comes from a learned function of the position, the bonus is no
longer calibrated to any confidence level, and it can be systematically wrong.
There is no published regret bound or consistency proof for PUCT with learned
priors. This is not a defect of the implementation; it is the state of the
theory, and Section~\ref{sec:proven} states it plainly.
\end{remark}

\paragraph{Proven values.} One classical refinement does carry over intact and is
used directly in this monograph. \citet{winands2008solver} observe that some
nodes have values that are \emph{exact}: a node whose game is over, or a node
whose value has been proven by a sound argument, does not need estimation. Their
MCTS-Solver propagates such values explicitly and stops sampling them. This is
the mechanism that lets a search recognise a forced win, and it is the same
mechanism that lets hand-written tactical detection feed exact labels into
training (Section~\ref{sec:alphaepsilon}).

\subsection{What is proven, and what is not}
\label{sec:proven}

The results collected in this section have very different status, and mixing them
is the most common way to overstate the strength of an AlphaZero-style system.
Table~\ref{tab:proven} separates them.

\begin{table}[t]
\centering
\small
\caption{Status of the results used by an AlphaZero-style agent. The last three
entries are the reason the systems work in practice; they are not theorems.}
\label{tab:proven}
\begin{tabularx}{\textwidth}{@{}X l@{}}
\toprule
Claim & Status \\
\midrule
UCB1 achieves logarithmic regret \citep{auer2002ucb1} & proven \\
UCT converges to the minimax value \citep{kocsis2006uct} & proven \\
Policy gradient with function approximation converges to a local
optimum \citep{sutton2000policygradient} & proven \\
The policy improvement theorem (greedy step does not harm the policy)
\citep{suttonbarto2018} & proven \\
Deterministic policy gradient formula \citep{silver2014dpg} & proven \\
Expert Iteration converges to the optimum \citep{anthony2017exit} & proven,
exact tabular case only \\
PUCT with learned priors converges & \textbf{not proven} \\
Approximate policy iteration with deep function approximation converges
& \textbf{not proven} \\
The AlphaZero loop (deep network $+$ MCTS $+$ self-play) converges
& \textbf{not proven} \\
\bottomrule
\end{tabularx}
\end{table}

Two consequences follow, and they shape the engineering in the rest of this
monograph. First, the guarantees that do exist are guarantees about the
\emph{skeleton} -- a bandit rule, an exact policy iteration, an exact expert.
Every component that makes the system practical breaks an assumption of those
proofs. Second, and consequently, correctness in an implementation must be
established by testing against known facts rather than by appeal to theory. In
the Gomoku setting, known facts are abundant: positions with exact values, and
forced-win sequences that can be verified mechanically. Section~\ref{sec:tree}
uses both.
